% new_system.tex -- design of the residency-arbitration runtime.
%
% STANDALONE-COMPILABLE so it can be read and checked on its own. To fold into
% the main paper, delete everything above \section{Design Overview} and
% everything below \end{document}, then % new_system.tex -- design of the residency-arbitration runtime.
%
% STANDALONE-COMPILABLE so it can be read and checked on its own. To fold into
% the main paper, delete everything above \section{Design Overview} and
% everything below \end{document}, then % new_system.tex -- design of the residency-arbitration runtime.
%
% STANDALONE-COMPILABLE so it can be read and checked on its own. To fold into
% the main paper, delete everything above \section{Design Overview} and
% everything below \end{document}, then % new_system.tex -- design of the residency-arbitration runtime.
%
% STANDALONE-COMPILABLE so it can be read and checked on its own. To fold into
% the main paper, delete everything above \section{Design Overview} and
% everything below \end{document}, then \input{new_system}.
%
% Packages used are all TeX Live / Overleaf defaults: geometry, booktabs,
% amsmath, xspace, algorithm, algpseudocode, hyperref.
%
% NOT COMPILED. No LaTeX toolchain was available on the machine this was
% written on, so treat a first build as a proofreading pass.
%
% PROVENANCE RULE, inherited from the plan and enforced throughout: every
% number carries "(measured)" or "(simulated)". Simulated numbers come from
% scripts/sim_residency_v2.py over synthetic schedules with measured resource
% constants; nothing labelled simulated ran on a GPU. A retention percentage
% is never quoted without the oracle-vs-LRU gap at the same budget.

\documentclass[10pt,twocolumn]{article}
\usepackage[margin=0.75in]{geometry}
\usepackage{booktabs}
\usepackage{amsmath}
\usepackage{xspace}
\usepackage{algorithm}
\usepackage{algpseudocode}
\usepackage[hidelinks]{hyperref}

\newcommand{\sysname}{Tandem\xspace}   % placeholder; rename in one place
\newcommand{\meas}{\textsuperscript{m}}
\newcommand{\simu}{\textsuperscript{s}}

\title{\sysname: Arbitrating Model and Data Residency\\Under One Memory Budget}
\date{}

\begin{document}
\maketitle

\noindent\textbf{Notation.} \meas{} marks a measured quantity;
\simu{} marks a simulation result over measured constants.

\section{Design Overview}

An agentic HPC workflow alternates between two kinds of resource: language
models that back the planner and its tool-calling steps, and scientific data
artifacts that the tools consume. Both are expensive to make usable and both
are used repeatedly, so both are candidates for retention. \sysname is a
runtime that decides, continuously and under a single host-memory budget,
which of them to keep resident and at which level.

The design rests on three claims, each measured before it was designed
around.

\paragraph{Claim 1: making a resource usable is a chain of costs, and
residency is a prefix of that chain already paid.}
We model residency as four rungs (Table~\ref{tab:ladder}). A scheduler's job
is to choose, per resource, the highest rung it can afford to hold.

\begin{table}[t]
\centering\small
\begin{tabular}{@{}llll@{}}
\toprule
Rung & State & Cost to reach & Kind \\
\midrule
R0 & on disk & --- & --- \\
R1 & bytes in page cache & disk $\to$ RAM & movement \\
R2 & bytes in a process & RAM $\to$ process & movement \\
R3 & activated in a live consumer & parse / construct & transformation \\
\bottomrule
\end{tabular}
\caption{The residency ladder. Byte-oriented tiering systems operate on
R0$\to$R1 exclusively; that is the definition of their interface, not a
limitation of their implementation.}
\label{tab:ladder}
\end{table}

\paragraph{Claim 2: the two resource classes sit at opposite ends of that
chain, so the correct action differs in kind.}
For a language model, \textbf{74--81\%\meas{}} of a cold boot is moving
weights and engine initialisation is $\approx$4.7\%\meas{}; the cost is
movement. For a data artifact the split inverts. Decomposing a 3.32\,GB
tabulated interatomic potential with per-process \texttt{getrusage} gives
2.0\%\meas{} disk$\to$RAM, 3.9\%\meas{} RAM$\to$process, and
\textbf{93.0\%\meas{}} parse and spline construction.

The consequence is not a tuning difference. Because movement dominates a
model, the winning action is to \emph{skip} it by never falling below R2, not
to schedule it earlier: parking eliminates both the disk$\to$RAM and
RAM$\to$GPU legs, whereas warming the page cache removes only the first and
changed measured boot time by $-7.0\%$ to $+37.1\%$\meas{} across eight nodes
--- sign-unstable. Because transformation dominates a data artifact,
pre-\emph{placing} its bytes recovers at most the 2.0\% movement share; the
only winning action is to have already run the transform, i.e. to hold the
artifact at R3 inside a live consumer.

\paragraph{Claim 3: the two savings live in the same place, which is what
makes this one system rather than two.}
A model parked at R2 and a database activated at R3 are both
\emph{anonymous memory inside a live process}. Neither is a file range, so
neither can be named, evicted, or repopulated by a byte-oriented tier; and
both draw on the same host-memory allocation. Two mechanisms with
incommensurable value functions, competing for one budget, is an arbitration
problem, and it is the contribution of this work.

\subsection{Architecture}

\sysname has four components:

\begin{enumerate}\itemsep2pt
\item \textbf{M1, model parking.} Park an outgoing inference engine at R2
      rather than terminating it.
\item \textbf{M2, activated-data retention.} Hold a parsed artifact at R3 in
      a resident worker rather than re-parsing per tool call.
\item \textbf{The arbitrator.} At every residency decision point, choose the
      set of resources to hold, across both classes, under one budget.
\item \textbf{The slack stager.} Opportunistically climb rungs for resources
      predicted to be needed soon, but only when doing so evicts nothing.
\end{enumerate}

\noindent A \emph{horizon predictor} supplies the single input the arbitrator
needs --- time until each resource is next required. Section~\ref{sec:pred}
explains why it is deliberately a secondary component.

\section{Mechanisms}

\subsection{M1: model parking}

M1 uses the inference server's existing sleep facility to move an engine's
weights from device memory into host memory while keeping the process, its
CUDA context, and its allocator alive. Measured on a 32B model at tp=2, a
cold boot of \textbf{782.27\,s} becomes a \textbf{2.076\,s} wake\meas{}, with
verbatim-identical generated text, unchanged \texttt{finish\_reason}, and
post-wake throughput of 20.05--20.08\,tok/s against a 20.20\,tok/s pre-sleep
steady state --- so no cost is deferred into generation. Wake is independent
of page-cache state (evicting 18 weight shards changed it by
0.015\,s\meas{}), which is the mechanistic confirmation that the weights are
held in anonymous memory rather than re-read from a file.

Parking is not free: it costs \textbf{1.90$\times$}\meas{} the weight bytes
in host RAM (120.77\,GiB to park a 68.28\,GB model). That ratio is what makes
residency a budgeted decision rather than a default.

\emph{We do not claim sleep/wake as a contribution.} It is a feature of the
serving stack, obtainable from a flag. What \sysname contributes on this axis
is deciding \emph{when} parking is affordable and \emph{what it should
displace} --- and, critically, that at realistic allocations the answer is
frequently ``not this model.''

\subsection{M2: activated-data retention}

Agent frameworks typically execute each tool call as a fresh process. The
activated structure therefore dies at every invocation, and the artifact is
re-parsed even though the bytes are unchanged and the page cache is already
warm --- a cost invisible to any byte-oriented prefetcher by the second
parse. M2 replaces fork-per-call with a resident worker that retains the
activated structure across calls: \textbf{93.73\,s $\to$ 10.56\,s}\meas{}
(8.88$\times$) per redundant invocation, with the physical output
bit-identical (relative difference $0.000\mathrm{e}{+}00$).

The mechanism is not free of semantics, and this is where most of the
engineering lives. A live simulation instance refuses to redefine its
simulation box, and the conventional reset command destroys precisely the
tables being retained. Reaching a new problem geometry while preserving them
requires an in-place empty--resize--refill sequence. Because \sysname
rewrites commands to achieve this, \textbf{a run that completes and returns
plausible numbers is not evidence of correctness}; retention is enabled only
for command shapes that have passed a differential test executing the same
script through both the retained and the fork-per-call path and comparing
physical outputs.

\section{Arbitration}

\subsection{One budget, two classes}

Let $\mathcal{R}$ be the resident-candidate set, $g(r)$ the host memory a
resource occupies while held, $B$ the budget, and $\Delta t(r)$ the predicted
time until $r$ is next needed. Define

\begin{align}
\mathrm{full}(r) &= \mathrm{cold}(r) - \mathrm{ready}(r) \\
\mathrm{benefit}(r) &=
  \begin{cases}
    \mathrm{full}(r) & \text{retain}\\
    \min\!\left(\mathrm{full}(r), \Delta t(r)\right) & \text{stage}
  \end{cases}\\
v(r) &= \mathrm{benefit}(r)\cdot \frac{H}{\max(\Delta t(r),\,H)}
\label{eq:value}
\end{align}

\noindent where $\mathrm{cold}$ is the cost of producing the resource from
R0, $\mathrm{ready}$ the cost of resuming it from its held rung, and $H$ a
planning horizon. The arbitrator solves
$\max \sum_{r\in S} v(r)$ subject to $\sum_{r \in S} g(r) \le B$.

\subsection{Why this value function and not the obvious one}

The natural choice --- benefit per byte per second, $v = \mathrm{benefit} /
\Delta t$ --- is wrong, and wrong in a way that dimensional analysis exposes
before any experiment. Retain and stage have benefit functions of different
shape in $\Delta t$: retain's is constant, stage's saturates at
$\min(\mathrm{full}, \Delta t)$. Dividing by $\Delta t$ therefore decays
retain's score as $1/\Delta t$ while leaving stage's at a flat $1/g(r)$
whenever the window is shorter than the load --- which is every interesting
case, gaps being seconds against loads of hundreds of seconds. The resulting
density has units of $\mathrm{GB}^{-1}$, not seconds per GB. Empirically, a
stage worth 5\,s outranks and evicts a retention worth 798\,s.

We test candidate objectives against an \textbf{oracle-monotonicity
invariant}: offering staging as an additional \emph{option} to an
oracle-informed policy must never lower the result, since the policy could
always decline. The rate objective violates it by $-2.05$ to $-3.77$
points\simu{}; the pure-benefit objective is marginal; the horizon objective
of Eq.~\ref{eq:value} passes at $+3.14$ to $+4.28$\simu{}. The invariant is
necessary, not sufficient --- all arms remain myopic one-step solves --- but
it is cheap and it caught a defect that had survived in an earlier
implementation.

$H$ is a real knob and is swept rather than chosen. Smaller $H$ is better
because the pathology lives at small $\Delta t$; any $H$ above the typical
inter-need gap suffices.

\subsection{Why the problem is 0/1 and not fractional}

Eviction here is a \emph{semantic} operation, not a paging one. Discarding a
model's weights is all-or-nothing, and terminating a worker destroys the
whole activated structure; there is no interface for parking half a model or
retaining 40\% of a database, and 90\% of a model's weights serve zero
tokens. Consequently the items are indivisible and the problem is 0/1
knapsack, for which density-ranked greedy is a heuristic rather than an
optimum. This is why \sysname evaluates a subset choice rather than
single-victim eviction: single-victim greedy cannot express ``evict these two
together to admit that one.'' Algorithm~\ref{alg:arb} states the decision
procedure; note that both retain and stage candidates are scored into the
\emph{same} objective, which is what makes the two actions comparable rather
than separately prioritised.

\begin{algorithm}[t]
\small
\caption{Arbitration at a residency decision point}
\label{alg:arb}
\begin{algorithmic}[1]
\State \textbf{input:} candidates $\mathcal{R}$, budget $B$, horizon oracle $\Delta t(\cdot)$
\State $\mathcal{A} \gets \{\,r \in \mathcal{R} : \Delta t(r) \ne \textsc{Never}\,\}$
  \Comment{believed-dead excluded, not scored zero}
\For{$r \in \mathcal{A}$}
  \State $w_r \gets v(r)$ via Eq.~\ref{eq:value} with kind \textsc{retain}
\EndFor
\For{$r \in \textsc{StageCandidates}$}
  \State $w_r \gets v(r)$ with kind \textsc{stage}
  \Comment{same currency as retain}
\EndFor
\State $S \gets \textsc{Knapsack}(w, g, B)$
  \Comment{ties break toward smaller $\sum g$}
\State \textbf{evict} $\mathrm{resident} \setminus S$; \textbf{admit} $S \setminus \mathrm{resident}$
\State \Return $S$
\end{algorithmic}
\end{algorithm}

\subsection{Interaction with device slots}

A model resident on an accelerator does not charge the host-memory budget,
since its weights are in device memory. The arbitrator therefore treats
device residency as a separate, smaller capacity and only considers for
host-memory retention what has been displaced from it. This has a
counter-intuitive consequence that the evaluation confirms: adding a device
slot \emph{reduces} the value of cost-aware host retention, because the
frequently-used models are already free and what remains competing for host
memory is the less-used tail.

\section{Speculative Staging Under Slack}

\sysname stages speculatively only under a strict safety rule: \textbf{a
stage may never cause an eviction}. If the predicted-next resource does not
fit in currently unused budget, the stage is abandoned rather than funded by
displacing a resident resource.

The rule is what makes staging safe at realistic prediction quality. Because
a stage can only ever consume memory nothing else wanted, its downside is
bounded by wasted bandwidth rather than by a lost retention, and its benefit
becomes nearly independent of predictor accuracy: $+3.83$ points at 0.40
accuracy versus $+5.06$ at a perfect oracle\simu{}. \emph{The safety rule
substitutes for accuracy.} This is the inverse of the usual prefetching
design point, where value scales with precision, and it is the direct
consequence of operating in a regime where the competing use of memory is
retention rather than nothing.

\section{The Horizon Predictor, and Why It Is Secondary}
\label{sec:pred}

\subsection{From identity to horizon}

An earlier version of this system used prediction the conventional way: infer
\emph{which} resource is needed next and fetch it. That framing failed, and
its failure motivated the present design. Two signals were available --- a
central planner's plan, and learned tool-to-tool transition statistics --- and
neither delivered usable identity prediction on non-deterministic agent
traces.

Under \sysname prediction answers a different question: \emph{how long until
resource $r$ is needed again?} It is queried per already-held resource at
eviction time, not consulted to select a fetch target. This is a strictly
weaker requirement, in a way that matters:

\begin{itemize}\itemsep2pt
\item A retention decision is made \emph{at eviction time}, when the past is
      known, rather than at fetch time, when it is not.
\item A wrong retain costs only memory. A wrong fetch costs bandwidth
      \emph{and} memory \emph{and} may leave the consumer stalled anyway.
\item Identity is not required at all: every held resource is scored, so
      there is no candidate to name.
\end{itemize}

\subsection{The accuracy the design actually needs}

Because the arbitrator ranks rather than selects, horizon error degrades the
ranking gracefully instead of committing to a wrong target. Across a sweep of
predictor accuracy from 0.15 to 1.00, the deployable density-ranked policy
moves from 47.63\% to 64.86\%\simu{} while a size-blind LRU baseline is flat
at 56.23\%\simu{} by construction. Break-even against LRU falls near
\textbf{0.55} --- the middle of the accuracy band the underlying predictor
actually achieves.

We state the negative half plainly: \emph{below that break-even, cost-aware
retention is worse than LRU}, and the entire oracle ceiling over LRU is about
nine points. The predictor is secondary not because it is unimportant but
because the mechanism captures most of the available benefit without it, and
the policy's marginal contribution is real but bounded.

\subsection{One non-obvious interface requirement}

A predictor must never report ``never again.'' In a finite trace, ``no
further use'' is an artifact of the trace ending, but the arbitrator treats it
as licence to discard. A bounded-lookahead predictor should instead report
\emph{not within $H$}, which preserves a resource's eligibility when budget
is available. Making this substitution is worth roughly nine points\simu{} to
the subset-selection arm at measured accuracy, and it converts an apparent
finding --- that exact optimisation amplifies prediction error --- into an
artifact of a dead/alive binary. Table~\ref{tab:iface} states the full
interface.

\begin{table}[t]
\centering\small
\begin{tabular}{@{}p{0.30\columnwidth}p{0.60\columnwidth}@{}}
\toprule
Requirement & Why \\
\midrule
Horizon in \emph{seconds} & The budget is a rate-vs-time trade; step counts
  do not compose with load costs. \\
Queried per \emph{held} resource & The decision is which of $n$ residents to
  drop, not which one thing to fetch. \\
Calibrated ``not within $H$'' & Never $\infty$; see above. \\
Queryable \emph{at eviction time} & The decision point is displacement, not
  tool dispatch. \\
Held footprint, not file size & The budget is charged the parked/activated
  size, e.g. $1.90\times$ weights. \\
Cold and ready costs, separately & $v(\cdot)$ needs
  $\mathrm{cold}-\mathrm{ready}$, not one load estimate. \\
Per-resource confidence & Used as a weight, so a low-confidence horizon
  degrades a rank rather than forcing a choice. \\
\bottomrule
\end{tabular}
\caption{What the horizon predictor must expose.}
\label{tab:iface}
\end{table}

\section{Scope Conditions}

The design is not universally profitable, and the conditions are sharp enough
to state rather than let a reader discover.

\paragraph{It requires a stall-dominated workload.} The benefit is a
fraction of time spent waiting for resources; where compute windows are large
the same seconds are a small share of wall time. This is defensible for the
workloads measured here --- a representative trial spends 280.3\,s of
5288.5\,s\meas{} doing agent work --- but it is a real boundary, and it runs
opposite to what speculative staging wants. \textbf{The two halves of the
system prefer opposite regimes}, which we report rather than reconcile.

\paragraph{It requires contention.} If the budget accommodates every
resource, the correct policy is ``hold everything'' and no arbitration is
needed. Conversely, if the largest resources exceed the budget alone they are
not badly ranked but unrankable --- at a 256\,GB allocation the two largest
models are individually 279.0 and 276.3\,GB\meas{}, so 51.4\%\simu{} of stall
is capacity rather than replacement, and no ranking can convert it.
Quantisation, not a better policy, is what moves that mass.

\paragraph{The value of each mechanism depends on the budget, and must be
quoted with it.} At a budget where nearly everything fits, retention is worth
a great deal and an oracle recovers essentially nothing over LRU --- the win
belongs to the serving stack's flag, not to arbitration. At the production
allocation the ordering inverts, and retaining \emph{both} classes under LRU
is worse than retaining data alone by 4.73 points\simu{}, because a
size-blind policy parks a large, rarely-used model over a smaller,
frequently-used artifact; cost-aware arbitration recovers $+1.29$\simu{}.
\textbf{Any retention figure reported without the oracle-vs-LRU gap at the
same budget is not a claim about this system.}

\paragraph{Transformation cost is format-determined, and formats can be
changed.} Seconds per GB retained spans $65\times$ across
formats\meas{} and is flat to $1.00$--$1.16\times$ across a $4\times$ size
range\meas{}, which is what makes it a cacheable scheduling input rather than
an online-learning problem. But an artifact whose consumer accepts a
pre-built binary form can have its transformation cost converted away
offline. That conversion is not free: it trades a low-variance CPU cost for a
larger byte-movement cost, and for the sequence database measured here the
two balance at \textbf{151\,MB/s}\meas{} --- above measured shared-filesystem
throughput, where pre-building is a net loss, and below local NVMe, where it
is roughly break-even. The honest position is that this is
configuration-dependent rather than decided.

\end{document}
.
%
% Packages used are all TeX Live / Overleaf defaults: geometry, booktabs,
% amsmath, xspace, algorithm, algpseudocode, hyperref.
%
% NOT COMPILED. No LaTeX toolchain was available on the machine this was
% written on, so treat a first build as a proofreading pass.
%
% PROVENANCE RULE, inherited from the plan and enforced throughout: every
% number carries "(measured)" or "(simulated)". Simulated numbers come from
% scripts/sim_residency_v2.py over synthetic schedules with measured resource
% constants; nothing labelled simulated ran on a GPU. A retention percentage
% is never quoted without the oracle-vs-LRU gap at the same budget.

\documentclass[10pt,twocolumn]{article}
\usepackage[margin=0.75in]{geometry}
\usepackage{booktabs}
\usepackage{amsmath}
\usepackage{xspace}
\usepackage{algorithm}
\usepackage{algpseudocode}
\usepackage[hidelinks]{hyperref}

\newcommand{\sysname}{Tandem\xspace}   % placeholder; rename in one place
\newcommand{\meas}{\textsuperscript{m}}
\newcommand{\simu}{\textsuperscript{s}}

\title{\sysname: Arbitrating Model and Data Residency\\Under One Memory Budget}
\date{}

\begin{document}
\maketitle

\noindent\textbf{Notation.} \meas{} marks a measured quantity;
\simu{} marks a simulation result over measured constants.

\section{Design Overview}

An agentic HPC workflow alternates between two kinds of resource: language
models that back the planner and its tool-calling steps, and scientific data
artifacts that the tools consume. Both are expensive to make usable and both
are used repeatedly, so both are candidates for retention. \sysname is a
runtime that decides, continuously and under a single host-memory budget,
which of them to keep resident and at which level.

The design rests on three claims, each measured before it was designed
around.

\paragraph{Claim 1: making a resource usable is a chain of costs, and
residency is a prefix of that chain already paid.}
We model residency as four rungs (Table~\ref{tab:ladder}). A scheduler's job
is to choose, per resource, the highest rung it can afford to hold.

\begin{table}[t]
\centering\small
\begin{tabular}{@{}llll@{}}
\toprule
Rung & State & Cost to reach & Kind \\
\midrule
R0 & on disk & --- & --- \\
R1 & bytes in page cache & disk $\to$ RAM & movement \\
R2 & bytes in a process & RAM $\to$ process & movement \\
R3 & activated in a live consumer & parse / construct & transformation \\
\bottomrule
\end{tabular}
\caption{The residency ladder. Byte-oriented tiering systems operate on
R0$\to$R1 exclusively; that is the definition of their interface, not a
limitation of their implementation.}
\label{tab:ladder}
\end{table}

\paragraph{Claim 2: the two resource classes sit at opposite ends of that
chain, so the correct action differs in kind.}
For a language model, \textbf{74--81\%\meas{}} of a cold boot is moving
weights and engine initialisation is $\approx$4.7\%\meas{}; the cost is
movement. For a data artifact the split inverts. Decomposing a 3.32\,GB
tabulated interatomic potential with per-process \texttt{getrusage} gives
2.0\%\meas{} disk$\to$RAM, 3.9\%\meas{} RAM$\to$process, and
\textbf{93.0\%\meas{}} parse and spline construction.

The consequence is not a tuning difference. Because movement dominates a
model, the winning action is to \emph{skip} it by never falling below R2, not
to schedule it earlier: parking eliminates both the disk$\to$RAM and
RAM$\to$GPU legs, whereas warming the page cache removes only the first and
changed measured boot time by $-7.0\%$ to $+37.1\%$\meas{} across eight nodes
--- sign-unstable. Because transformation dominates a data artifact,
pre-\emph{placing} its bytes recovers at most the 2.0\% movement share; the
only winning action is to have already run the transform, i.e. to hold the
artifact at R3 inside a live consumer.

\paragraph{Claim 3: the two savings live in the same place, which is what
makes this one system rather than two.}
A model parked at R2 and a database activated at R3 are both
\emph{anonymous memory inside a live process}. Neither is a file range, so
neither can be named, evicted, or repopulated by a byte-oriented tier; and
both draw on the same host-memory allocation. Two mechanisms with
incommensurable value functions, competing for one budget, is an arbitration
problem, and it is the contribution of this work.

\subsection{Architecture}

\sysname has four components:

\begin{enumerate}\itemsep2pt
\item \textbf{M1, model parking.} Park an outgoing inference engine at R2
      rather than terminating it.
\item \textbf{M2, activated-data retention.} Hold a parsed artifact at R3 in
      a resident worker rather than re-parsing per tool call.
\item \textbf{The arbitrator.} At every residency decision point, choose the
      set of resources to hold, across both classes, under one budget.
\item \textbf{The slack stager.} Opportunistically climb rungs for resources
      predicted to be needed soon, but only when doing so evicts nothing.
\end{enumerate}

\noindent A \emph{horizon predictor} supplies the single input the arbitrator
needs --- time until each resource is next required. Section~\ref{sec:pred}
explains why it is deliberately a secondary component.

\section{Mechanisms}

\subsection{M1: model parking}

M1 uses the inference server's existing sleep facility to move an engine's
weights from device memory into host memory while keeping the process, its
CUDA context, and its allocator alive. Measured on a 32B model at tp=2, a
cold boot of \textbf{782.27\,s} becomes a \textbf{2.076\,s} wake\meas{}, with
verbatim-identical generated text, unchanged \texttt{finish\_reason}, and
post-wake throughput of 20.05--20.08\,tok/s against a 20.20\,tok/s pre-sleep
steady state --- so no cost is deferred into generation. Wake is independent
of page-cache state (evicting 18 weight shards changed it by
0.015\,s\meas{}), which is the mechanistic confirmation that the weights are
held in anonymous memory rather than re-read from a file.

Parking is not free: it costs \textbf{1.90$\times$}\meas{} the weight bytes
in host RAM (120.77\,GiB to park a 68.28\,GB model). That ratio is what makes
residency a budgeted decision rather than a default.

\emph{We do not claim sleep/wake as a contribution.} It is a feature of the
serving stack, obtainable from a flag. What \sysname contributes on this axis
is deciding \emph{when} parking is affordable and \emph{what it should
displace} --- and, critically, that at realistic allocations the answer is
frequently ``not this model.''

\subsection{M2: activated-data retention}

Agent frameworks typically execute each tool call as a fresh process. The
activated structure therefore dies at every invocation, and the artifact is
re-parsed even though the bytes are unchanged and the page cache is already
warm --- a cost invisible to any byte-oriented prefetcher by the second
parse. M2 replaces fork-per-call with a resident worker that retains the
activated structure across calls: \textbf{93.73\,s $\to$ 10.56\,s}\meas{}
(8.88$\times$) per redundant invocation, with the physical output
bit-identical (relative difference $0.000\mathrm{e}{+}00$).

The mechanism is not free of semantics, and this is where most of the
engineering lives. A live simulation instance refuses to redefine its
simulation box, and the conventional reset command destroys precisely the
tables being retained. Reaching a new problem geometry while preserving them
requires an in-place empty--resize--refill sequence. Because \sysname
rewrites commands to achieve this, \textbf{a run that completes and returns
plausible numbers is not evidence of correctness}; retention is enabled only
for command shapes that have passed a differential test executing the same
script through both the retained and the fork-per-call path and comparing
physical outputs.

\section{Arbitration}

\subsection{One budget, two classes}

Let $\mathcal{R}$ be the resident-candidate set, $g(r)$ the host memory a
resource occupies while held, $B$ the budget, and $\Delta t(r)$ the predicted
time until $r$ is next needed. Define

\begin{align}
\mathrm{full}(r) &= \mathrm{cold}(r) - \mathrm{ready}(r) \\
\mathrm{benefit}(r) &=
  \begin{cases}
    \mathrm{full}(r) & \text{retain}\\
    \min\!\left(\mathrm{full}(r), \Delta t(r)\right) & \text{stage}
  \end{cases}\\
v(r) &= \mathrm{benefit}(r)\cdot \frac{H}{\max(\Delta t(r),\,H)}
\label{eq:value}
\end{align}

\noindent where $\mathrm{cold}$ is the cost of producing the resource from
R0, $\mathrm{ready}$ the cost of resuming it from its held rung, and $H$ a
planning horizon. The arbitrator solves
$\max \sum_{r\in S} v(r)$ subject to $\sum_{r \in S} g(r) \le B$.

\subsection{Why this value function and not the obvious one}

The natural choice --- benefit per byte per second, $v = \mathrm{benefit} /
\Delta t$ --- is wrong, and wrong in a way that dimensional analysis exposes
before any experiment. Retain and stage have benefit functions of different
shape in $\Delta t$: retain's is constant, stage's saturates at
$\min(\mathrm{full}, \Delta t)$. Dividing by $\Delta t$ therefore decays
retain's score as $1/\Delta t$ while leaving stage's at a flat $1/g(r)$
whenever the window is shorter than the load --- which is every interesting
case, gaps being seconds against loads of hundreds of seconds. The resulting
density has units of $\mathrm{GB}^{-1}$, not seconds per GB. Empirically, a
stage worth 5\,s outranks and evicts a retention worth 798\,s.

We test candidate objectives against an \textbf{oracle-monotonicity
invariant}: offering staging as an additional \emph{option} to an
oracle-informed policy must never lower the result, since the policy could
always decline. The rate objective violates it by $-2.05$ to $-3.77$
points\simu{}; the pure-benefit objective is marginal; the horizon objective
of Eq.~\ref{eq:value} passes at $+3.14$ to $+4.28$\simu{}. The invariant is
necessary, not sufficient --- all arms remain myopic one-step solves --- but
it is cheap and it caught a defect that had survived in an earlier
implementation.

$H$ is a real knob and is swept rather than chosen. Smaller $H$ is better
because the pathology lives at small $\Delta t$; any $H$ above the typical
inter-need gap suffices.

\subsection{Why the problem is 0/1 and not fractional}

Eviction here is a \emph{semantic} operation, not a paging one. Discarding a
model's weights is all-or-nothing, and terminating a worker destroys the
whole activated structure; there is no interface for parking half a model or
retaining 40\% of a database, and 90\% of a model's weights serve zero
tokens. Consequently the items are indivisible and the problem is 0/1
knapsack, for which density-ranked greedy is a heuristic rather than an
optimum. This is why \sysname evaluates a subset choice rather than
single-victim eviction: single-victim greedy cannot express ``evict these two
together to admit that one.'' Algorithm~\ref{alg:arb} states the decision
procedure; note that both retain and stage candidates are scored into the
\emph{same} objective, which is what makes the two actions comparable rather
than separately prioritised.

\begin{algorithm}[t]
\small
\caption{Arbitration at a residency decision point}
\label{alg:arb}
\begin{algorithmic}[1]
\State \textbf{input:} candidates $\mathcal{R}$, budget $B$, horizon oracle $\Delta t(\cdot)$
\State $\mathcal{A} \gets \{\,r \in \mathcal{R} : \Delta t(r) \ne \textsc{Never}\,\}$
  \Comment{believed-dead excluded, not scored zero}
\For{$r \in \mathcal{A}$}
  \State $w_r \gets v(r)$ via Eq.~\ref{eq:value} with kind \textsc{retain}
\EndFor
\For{$r \in \textsc{StageCandidates}$}
  \State $w_r \gets v(r)$ with kind \textsc{stage}
  \Comment{same currency as retain}
\EndFor
\State $S \gets \textsc{Knapsack}(w, g, B)$
  \Comment{ties break toward smaller $\sum g$}
\State \textbf{evict} $\mathrm{resident} \setminus S$; \textbf{admit} $S \setminus \mathrm{resident}$
\State \Return $S$
\end{algorithmic}
\end{algorithm}

\subsection{Interaction with device slots}

A model resident on an accelerator does not charge the host-memory budget,
since its weights are in device memory. The arbitrator therefore treats
device residency as a separate, smaller capacity and only considers for
host-memory retention what has been displaced from it. This has a
counter-intuitive consequence that the evaluation confirms: adding a device
slot \emph{reduces} the value of cost-aware host retention, because the
frequently-used models are already free and what remains competing for host
memory is the less-used tail.

\section{Speculative Staging Under Slack}

\sysname stages speculatively only under a strict safety rule: \textbf{a
stage may never cause an eviction}. If the predicted-next resource does not
fit in currently unused budget, the stage is abandoned rather than funded by
displacing a resident resource.

The rule is what makes staging safe at realistic prediction quality. Because
a stage can only ever consume memory nothing else wanted, its downside is
bounded by wasted bandwidth rather than by a lost retention, and its benefit
becomes nearly independent of predictor accuracy: $+3.83$ points at 0.40
accuracy versus $+5.06$ at a perfect oracle\simu{}. \emph{The safety rule
substitutes for accuracy.} This is the inverse of the usual prefetching
design point, where value scales with precision, and it is the direct
consequence of operating in a regime where the competing use of memory is
retention rather than nothing.

\section{The Horizon Predictor, and Why It Is Secondary}
\label{sec:pred}

\subsection{From identity to horizon}

An earlier version of this system used prediction the conventional way: infer
\emph{which} resource is needed next and fetch it. That framing failed, and
its failure motivated the present design. Two signals were available --- a
central planner's plan, and learned tool-to-tool transition statistics --- and
neither delivered usable identity prediction on non-deterministic agent
traces.

Under \sysname prediction answers a different question: \emph{how long until
resource $r$ is needed again?} It is queried per already-held resource at
eviction time, not consulted to select a fetch target. This is a strictly
weaker requirement, in a way that matters:

\begin{itemize}\itemsep2pt
\item A retention decision is made \emph{at eviction time}, when the past is
      known, rather than at fetch time, when it is not.
\item A wrong retain costs only memory. A wrong fetch costs bandwidth
      \emph{and} memory \emph{and} may leave the consumer stalled anyway.
\item Identity is not required at all: every held resource is scored, so
      there is no candidate to name.
\end{itemize}

\subsection{The accuracy the design actually needs}

Because the arbitrator ranks rather than selects, horizon error degrades the
ranking gracefully instead of committing to a wrong target. Across a sweep of
predictor accuracy from 0.15 to 1.00, the deployable density-ranked policy
moves from 47.63\% to 64.86\%\simu{} while a size-blind LRU baseline is flat
at 56.23\%\simu{} by construction. Break-even against LRU falls near
\textbf{0.55} --- the middle of the accuracy band the underlying predictor
actually achieves.

We state the negative half plainly: \emph{below that break-even, cost-aware
retention is worse than LRU}, and the entire oracle ceiling over LRU is about
nine points. The predictor is secondary not because it is unimportant but
because the mechanism captures most of the available benefit without it, and
the policy's marginal contribution is real but bounded.

\subsection{One non-obvious interface requirement}

A predictor must never report ``never again.'' In a finite trace, ``no
further use'' is an artifact of the trace ending, but the arbitrator treats it
as licence to discard. A bounded-lookahead predictor should instead report
\emph{not within $H$}, which preserves a resource's eligibility when budget
is available. Making this substitution is worth roughly nine points\simu{} to
the subset-selection arm at measured accuracy, and it converts an apparent
finding --- that exact optimisation amplifies prediction error --- into an
artifact of a dead/alive binary. Table~\ref{tab:iface} states the full
interface.

\begin{table}[t]
\centering\small
\begin{tabular}{@{}p{0.30\columnwidth}p{0.60\columnwidth}@{}}
\toprule
Requirement & Why \\
\midrule
Horizon in \emph{seconds} & The budget is a rate-vs-time trade; step counts
  do not compose with load costs. \\
Queried per \emph{held} resource & The decision is which of $n$ residents to
  drop, not which one thing to fetch. \\
Calibrated ``not within $H$'' & Never $\infty$; see above. \\
Queryable \emph{at eviction time} & The decision point is displacement, not
  tool dispatch. \\
Held footprint, not file size & The budget is charged the parked/activated
  size, e.g. $1.90\times$ weights. \\
Cold and ready costs, separately & $v(\cdot)$ needs
  $\mathrm{cold}-\mathrm{ready}$, not one load estimate. \\
Per-resource confidence & Used as a weight, so a low-confidence horizon
  degrades a rank rather than forcing a choice. \\
\bottomrule
\end{tabular}
\caption{What the horizon predictor must expose.}
\label{tab:iface}
\end{table}

\section{Scope Conditions}

The design is not universally profitable, and the conditions are sharp enough
to state rather than let a reader discover.

\paragraph{It requires a stall-dominated workload.} The benefit is a
fraction of time spent waiting for resources; where compute windows are large
the same seconds are a small share of wall time. This is defensible for the
workloads measured here --- a representative trial spends 280.3\,s of
5288.5\,s\meas{} doing agent work --- but it is a real boundary, and it runs
opposite to what speculative staging wants. \textbf{The two halves of the
system prefer opposite regimes}, which we report rather than reconcile.

\paragraph{It requires contention.} If the budget accommodates every
resource, the correct policy is ``hold everything'' and no arbitration is
needed. Conversely, if the largest resources exceed the budget alone they are
not badly ranked but unrankable --- at a 256\,GB allocation the two largest
models are individually 279.0 and 276.3\,GB\meas{}, so 51.4\%\simu{} of stall
is capacity rather than replacement, and no ranking can convert it.
Quantisation, not a better policy, is what moves that mass.

\paragraph{The value of each mechanism depends on the budget, and must be
quoted with it.} At a budget where nearly everything fits, retention is worth
a great deal and an oracle recovers essentially nothing over LRU --- the win
belongs to the serving stack's flag, not to arbitration. At the production
allocation the ordering inverts, and retaining \emph{both} classes under LRU
is worse than retaining data alone by 4.73 points\simu{}, because a
size-blind policy parks a large, rarely-used model over a smaller,
frequently-used artifact; cost-aware arbitration recovers $+1.29$\simu{}.
\textbf{Any retention figure reported without the oracle-vs-LRU gap at the
same budget is not a claim about this system.}

\paragraph{Transformation cost is format-determined, and formats can be
changed.} Seconds per GB retained spans $65\times$ across
formats\meas{} and is flat to $1.00$--$1.16\times$ across a $4\times$ size
range\meas{}, which is what makes it a cacheable scheduling input rather than
an online-learning problem. But an artifact whose consumer accepts a
pre-built binary form can have its transformation cost converted away
offline. That conversion is not free: it trades a low-variance CPU cost for a
larger byte-movement cost, and for the sequence database measured here the
two balance at \textbf{151\,MB/s}\meas{} --- above measured shared-filesystem
throughput, where pre-building is a net loss, and below local NVMe, where it
is roughly break-even. The honest position is that this is
configuration-dependent rather than decided.

\end{document}
.
%
% Packages used are all TeX Live / Overleaf defaults: geometry, booktabs,
% amsmath, xspace, algorithm, algpseudocode, hyperref.
%
% NOT COMPILED. No LaTeX toolchain was available on the machine this was
% written on, so treat a first build as a proofreading pass.
%
% PROVENANCE RULE, inherited from the plan and enforced throughout: every
% number carries "(measured)" or "(simulated)". Simulated numbers come from
% scripts/sim_residency_v2.py over synthetic schedules with measured resource
% constants; nothing labelled simulated ran on a GPU. A retention percentage
% is never quoted without the oracle-vs-LRU gap at the same budget.

\documentclass[10pt,twocolumn]{article}
\usepackage[margin=0.75in]{geometry}
\usepackage{booktabs}
\usepackage{amsmath}
\usepackage{xspace}
\usepackage{algorithm}
\usepackage{algpseudocode}
\usepackage[hidelinks]{hyperref}

\newcommand{\sysname}{Tandem\xspace}   % placeholder; rename in one place
\newcommand{\meas}{\textsuperscript{m}}
\newcommand{\simu}{\textsuperscript{s}}

\title{\sysname: Arbitrating Model and Data Residency\\Under One Memory Budget}
\date{}

\begin{document}
\maketitle

\noindent\textbf{Notation.} \meas{} marks a measured quantity;
\simu{} marks a simulation result over measured constants.

\section{Design Overview}

An agentic HPC workflow alternates between two kinds of resource: language
models that back the planner and its tool-calling steps, and scientific data
artifacts that the tools consume. Both are expensive to make usable and both
are used repeatedly, so both are candidates for retention. \sysname is a
runtime that decides, continuously and under a single host-memory budget,
which of them to keep resident and at which level.

The design rests on three claims, each measured before it was designed
around.

\paragraph{Claim 1: making a resource usable is a chain of costs, and
residency is a prefix of that chain already paid.}
We model residency as four rungs (Table~\ref{tab:ladder}). A scheduler's job
is to choose, per resource, the highest rung it can afford to hold.

\begin{table}[t]
\centering\small
\begin{tabular}{@{}llll@{}}
\toprule
Rung & State & Cost to reach & Kind \\
\midrule
R0 & on disk & --- & --- \\
R1 & bytes in page cache & disk $\to$ RAM & movement \\
R2 & bytes in a process & RAM $\to$ process & movement \\
R3 & activated in a live consumer & parse / construct & transformation \\
\bottomrule
\end{tabular}
\caption{The residency ladder. Byte-oriented tiering systems operate on
R0$\to$R1 exclusively; that is the definition of their interface, not a
limitation of their implementation.}
\label{tab:ladder}
\end{table}

\paragraph{Claim 2: the two resource classes sit at opposite ends of that
chain, so the correct action differs in kind.}
For a language model, \textbf{74--81\%\meas{}} of a cold boot is moving
weights and engine initialisation is $\approx$4.7\%\meas{}; the cost is
movement. For a data artifact the split inverts. Decomposing a 3.32\,GB
tabulated interatomic potential with per-process \texttt{getrusage} gives
2.0\%\meas{} disk$\to$RAM, 3.9\%\meas{} RAM$\to$process, and
\textbf{93.0\%\meas{}} parse and spline construction.

The consequence is not a tuning difference. Because movement dominates a
model, the winning action is to \emph{skip} it by never falling below R2, not
to schedule it earlier: parking eliminates both the disk$\to$RAM and
RAM$\to$GPU legs, whereas warming the page cache removes only the first and
changed measured boot time by $-7.0\%$ to $+37.1\%$\meas{} across eight nodes
--- sign-unstable. Because transformation dominates a data artifact,
pre-\emph{placing} its bytes recovers at most the 2.0\% movement share; the
only winning action is to have already run the transform, i.e. to hold the
artifact at R3 inside a live consumer.

\paragraph{Claim 3: the two savings live in the same place, which is what
makes this one system rather than two.}
A model parked at R2 and a database activated at R3 are both
\emph{anonymous memory inside a live process}. Neither is a file range, so
neither can be named, evicted, or repopulated by a byte-oriented tier; and
both draw on the same host-memory allocation. Two mechanisms with
incommensurable value functions, competing for one budget, is an arbitration
problem, and it is the contribution of this work.

\subsection{Architecture}

\sysname has four components:

\begin{enumerate}\itemsep2pt
\item \textbf{M1, model parking.} Park an outgoing inference engine at R2
      rather than terminating it.
\item \textbf{M2, activated-data retention.} Hold a parsed artifact at R3 in
      a resident worker rather than re-parsing per tool call.
\item \textbf{The arbitrator.} At every residency decision point, choose the
      set of resources to hold, across both classes, under one budget.
\item \textbf{The slack stager.} Opportunistically climb rungs for resources
      predicted to be needed soon, but only when doing so evicts nothing.
\end{enumerate}

\noindent A \emph{horizon predictor} supplies the single input the arbitrator
needs --- time until each resource is next required. Section~\ref{sec:pred}
explains why it is deliberately a secondary component.

\section{Mechanisms}

\subsection{M1: model parking}

M1 uses the inference server's existing sleep facility to move an engine's
weights from device memory into host memory while keeping the process, its
CUDA context, and its allocator alive. Measured on a 32B model at tp=2, a
cold boot of \textbf{782.27\,s} becomes a \textbf{2.076\,s} wake\meas{}, with
verbatim-identical generated text, unchanged \texttt{finish\_reason}, and
post-wake throughput of 20.05--20.08\,tok/s against a 20.20\,tok/s pre-sleep
steady state --- so no cost is deferred into generation. Wake is independent
of page-cache state (evicting 18 weight shards changed it by
0.015\,s\meas{}), which is the mechanistic confirmation that the weights are
held in anonymous memory rather than re-read from a file.

Parking is not free: it costs \textbf{1.90$\times$}\meas{} the weight bytes
in host RAM (120.77\,GiB to park a 68.28\,GB model). That ratio is what makes
residency a budgeted decision rather than a default.

\emph{We do not claim sleep/wake as a contribution.} It is a feature of the
serving stack, obtainable from a flag. What \sysname contributes on this axis
is deciding \emph{when} parking is affordable and \emph{what it should
displace} --- and, critically, that at realistic allocations the answer is
frequently ``not this model.''

\subsection{M2: activated-data retention}

Agent frameworks typically execute each tool call as a fresh process. The
activated structure therefore dies at every invocation, and the artifact is
re-parsed even though the bytes are unchanged and the page cache is already
warm --- a cost invisible to any byte-oriented prefetcher by the second
parse. M2 replaces fork-per-call with a resident worker that retains the
activated structure across calls: \textbf{93.73\,s $\to$ 10.56\,s}\meas{}
(8.88$\times$) per redundant invocation, with the physical output
bit-identical (relative difference $0.000\mathrm{e}{+}00$).

The mechanism is not free of semantics, and this is where most of the
engineering lives. A live simulation instance refuses to redefine its
simulation box, and the conventional reset command destroys precisely the
tables being retained. Reaching a new problem geometry while preserving them
requires an in-place empty--resize--refill sequence. Because \sysname
rewrites commands to achieve this, \textbf{a run that completes and returns
plausible numbers is not evidence of correctness}; retention is enabled only
for command shapes that have passed a differential test executing the same
script through both the retained and the fork-per-call path and comparing
physical outputs.

\section{Arbitration}

\subsection{One budget, two classes}

Let $\mathcal{R}$ be the resident-candidate set, $g(r)$ the host memory a
resource occupies while held, $B$ the budget, and $\Delta t(r)$ the predicted
time until $r$ is next needed. Define

\begin{align}
\mathrm{full}(r) &= \mathrm{cold}(r) - \mathrm{ready}(r) \\
\mathrm{benefit}(r) &=
  \begin{cases}
    \mathrm{full}(r) & \text{retain}\\
    \min\!\left(\mathrm{full}(r), \Delta t(r)\right) & \text{stage}
  \end{cases}\\
v(r) &= \mathrm{benefit}(r)\cdot \frac{H}{\max(\Delta t(r),\,H)}
\label{eq:value}
\end{align}

\noindent where $\mathrm{cold}$ is the cost of producing the resource from
R0, $\mathrm{ready}$ the cost of resuming it from its held rung, and $H$ a
planning horizon. The arbitrator solves
$\max \sum_{r\in S} v(r)$ subject to $\sum_{r \in S} g(r) \le B$.

\subsection{Why this value function and not the obvious one}

The natural choice --- benefit per byte per second, $v = \mathrm{benefit} /
\Delta t$ --- is wrong, and wrong in a way that dimensional analysis exposes
before any experiment. Retain and stage have benefit functions of different
shape in $\Delta t$: retain's is constant, stage's saturates at
$\min(\mathrm{full}, \Delta t)$. Dividing by $\Delta t$ therefore decays
retain's score as $1/\Delta t$ while leaving stage's at a flat $1/g(r)$
whenever the window is shorter than the load --- which is every interesting
case, gaps being seconds against loads of hundreds of seconds. The resulting
density has units of $\mathrm{GB}^{-1}$, not seconds per GB. Empirically, a
stage worth 5\,s outranks and evicts a retention worth 798\,s.

We test candidate objectives against an \textbf{oracle-monotonicity
invariant}: offering staging as an additional \emph{option} to an
oracle-informed policy must never lower the result, since the policy could
always decline. The rate objective violates it by $-2.05$ to $-3.77$
points\simu{}; the pure-benefit objective is marginal; the horizon objective
of Eq.~\ref{eq:value} passes at $+3.14$ to $+4.28$\simu{}. The invariant is
necessary, not sufficient --- all arms remain myopic one-step solves --- but
it is cheap and it caught a defect that had survived in an earlier
implementation.

$H$ is a real knob and is swept rather than chosen. Smaller $H$ is better
because the pathology lives at small $\Delta t$; any $H$ above the typical
inter-need gap suffices.

\subsection{Why the problem is 0/1 and not fractional}

Eviction here is a \emph{semantic} operation, not a paging one. Discarding a
model's weights is all-or-nothing, and terminating a worker destroys the
whole activated structure; there is no interface for parking half a model or
retaining 40\% of a database, and 90\% of a model's weights serve zero
tokens. Consequently the items are indivisible and the problem is 0/1
knapsack, for which density-ranked greedy is a heuristic rather than an
optimum. This is why \sysname evaluates a subset choice rather than
single-victim eviction: single-victim greedy cannot express ``evict these two
together to admit that one.'' Algorithm~\ref{alg:arb} states the decision
procedure; note that both retain and stage candidates are scored into the
\emph{same} objective, which is what makes the two actions comparable rather
than separately prioritised.

\begin{algorithm}[t]
\small
\caption{Arbitration at a residency decision point}
\label{alg:arb}
\begin{algorithmic}[1]
\State \textbf{input:} candidates $\mathcal{R}$, budget $B$, horizon oracle $\Delta t(\cdot)$
\State $\mathcal{A} \gets \{\,r \in \mathcal{R} : \Delta t(r) \ne \textsc{Never}\,\}$
  \Comment{believed-dead excluded, not scored zero}
\For{$r \in \mathcal{A}$}
  \State $w_r \gets v(r)$ via Eq.~\ref{eq:value} with kind \textsc{retain}
\EndFor
\For{$r \in \textsc{StageCandidates}$}
  \State $w_r \gets v(r)$ with kind \textsc{stage}
  \Comment{same currency as retain}
\EndFor
\State $S \gets \textsc{Knapsack}(w, g, B)$
  \Comment{ties break toward smaller $\sum g$}
\State \textbf{evict} $\mathrm{resident} \setminus S$; \textbf{admit} $S \setminus \mathrm{resident}$
\State \Return $S$
\end{algorithmic}
\end{algorithm}

\subsection{Interaction with device slots}

A model resident on an accelerator does not charge the host-memory budget,
since its weights are in device memory. The arbitrator therefore treats
device residency as a separate, smaller capacity and only considers for
host-memory retention what has been displaced from it. This has a
counter-intuitive consequence that the evaluation confirms: adding a device
slot \emph{reduces} the value of cost-aware host retention, because the
frequently-used models are already free and what remains competing for host
memory is the less-used tail.

\section{Speculative Staging Under Slack}

\sysname stages speculatively only under a strict safety rule: \textbf{a
stage may never cause an eviction}. If the predicted-next resource does not
fit in currently unused budget, the stage is abandoned rather than funded by
displacing a resident resource.

The rule is what makes staging safe at realistic prediction quality. Because
a stage can only ever consume memory nothing else wanted, its downside is
bounded by wasted bandwidth rather than by a lost retention, and its benefit
becomes nearly independent of predictor accuracy: $+3.83$ points at 0.40
accuracy versus $+5.06$ at a perfect oracle\simu{}. \emph{The safety rule
substitutes for accuracy.} This is the inverse of the usual prefetching
design point, where value scales with precision, and it is the direct
consequence of operating in a regime where the competing use of memory is
retention rather than nothing.

\section{The Horizon Predictor, and Why It Is Secondary}
\label{sec:pred}

\subsection{From identity to horizon}

An earlier version of this system used prediction the conventional way: infer
\emph{which} resource is needed next and fetch it. That framing failed, and
its failure motivated the present design. Two signals were available --- a
central planner's plan, and learned tool-to-tool transition statistics --- and
neither delivered usable identity prediction on non-deterministic agent
traces.

Under \sysname prediction answers a different question: \emph{how long until
resource $r$ is needed again?} It is queried per already-held resource at
eviction time, not consulted to select a fetch target. This is a strictly
weaker requirement, in a way that matters:

\begin{itemize}\itemsep2pt
\item A retention decision is made \emph{at eviction time}, when the past is
      known, rather than at fetch time, when it is not.
\item A wrong retain costs only memory. A wrong fetch costs bandwidth
      \emph{and} memory \emph{and} may leave the consumer stalled anyway.
\item Identity is not required at all: every held resource is scored, so
      there is no candidate to name.
\end{itemize}

\subsection{The accuracy the design actually needs}

Because the arbitrator ranks rather than selects, horizon error degrades the
ranking gracefully instead of committing to a wrong target. Across a sweep of
predictor accuracy from 0.15 to 1.00, the deployable density-ranked policy
moves from 47.63\% to 64.86\%\simu{} while a size-blind LRU baseline is flat
at 56.23\%\simu{} by construction. Break-even against LRU falls near
\textbf{0.55} --- the middle of the accuracy band the underlying predictor
actually achieves.

We state the negative half plainly: \emph{below that break-even, cost-aware
retention is worse than LRU}, and the entire oracle ceiling over LRU is about
nine points. The predictor is secondary not because it is unimportant but
because the mechanism captures most of the available benefit without it, and
the policy's marginal contribution is real but bounded.

\subsection{One non-obvious interface requirement}

A predictor must never report ``never again.'' In a finite trace, ``no
further use'' is an artifact of the trace ending, but the arbitrator treats it
as licence to discard. A bounded-lookahead predictor should instead report
\emph{not within $H$}, which preserves a resource's eligibility when budget
is available. Making this substitution is worth roughly nine points\simu{} to
the subset-selection arm at measured accuracy, and it converts an apparent
finding --- that exact optimisation amplifies prediction error --- into an
artifact of a dead/alive binary. Table~\ref{tab:iface} states the full
interface.

\begin{table}[t]
\centering\small
\begin{tabular}{@{}p{0.30\columnwidth}p{0.60\columnwidth}@{}}
\toprule
Requirement & Why \\
\midrule
Horizon in \emph{seconds} & The budget is a rate-vs-time trade; step counts
  do not compose with load costs. \\
Queried per \emph{held} resource & The decision is which of $n$ residents to
  drop, not which one thing to fetch. \\
Calibrated ``not within $H$'' & Never $\infty$; see above. \\
Queryable \emph{at eviction time} & The decision point is displacement, not
  tool dispatch. \\
Held footprint, not file size & The budget is charged the parked/activated
  size, e.g. $1.90\times$ weights. \\
Cold and ready costs, separately & $v(\cdot)$ needs
  $\mathrm{cold}-\mathrm{ready}$, not one load estimate. \\
Per-resource confidence & Used as a weight, so a low-confidence horizon
  degrades a rank rather than forcing a choice. \\
\bottomrule
\end{tabular}
\caption{What the horizon predictor must expose.}
\label{tab:iface}
\end{table}

\section{Scope Conditions}

The design is not universally profitable, and the conditions are sharp enough
to state rather than let a reader discover.

\paragraph{It requires a stall-dominated workload.} The benefit is a
fraction of time spent waiting for resources; where compute windows are large
the same seconds are a small share of wall time. This is defensible for the
workloads measured here --- a representative trial spends 280.3\,s of
5288.5\,s\meas{} doing agent work --- but it is a real boundary, and it runs
opposite to what speculative staging wants. \textbf{The two halves of the
system prefer opposite regimes}, which we report rather than reconcile.

\paragraph{It requires contention.} If the budget accommodates every
resource, the correct policy is ``hold everything'' and no arbitration is
needed. Conversely, if the largest resources exceed the budget alone they are
not badly ranked but unrankable --- at a 256\,GB allocation the two largest
models are individually 279.0 and 276.3\,GB\meas{}, so 51.4\%\simu{} of stall
is capacity rather than replacement, and no ranking can convert it.
Quantisation, not a better policy, is what moves that mass.

\paragraph{The value of each mechanism depends on the budget, and must be
quoted with it.} At a budget where nearly everything fits, retention is worth
a great deal and an oracle recovers essentially nothing over LRU --- the win
belongs to the serving stack's flag, not to arbitration. At the production
allocation the ordering inverts, and retaining \emph{both} classes under LRU
is worse than retaining data alone by 4.73 points\simu{}, because a
size-blind policy parks a large, rarely-used model over a smaller,
frequently-used artifact; cost-aware arbitration recovers $+1.29$\simu{}.
\textbf{Any retention figure reported without the oracle-vs-LRU gap at the
same budget is not a claim about this system.}

\paragraph{Transformation cost is format-determined, and formats can be
changed.} Seconds per GB retained spans $65\times$ across
formats\meas{} and is flat to $1.00$--$1.16\times$ across a $4\times$ size
range\meas{}, which is what makes it a cacheable scheduling input rather than
an online-learning problem. But an artifact whose consumer accepts a
pre-built binary form can have its transformation cost converted away
offline. That conversion is not free: it trades a low-variance CPU cost for a
larger byte-movement cost, and for the sequence database measured here the
two balance at \textbf{151\,MB/s}\meas{} --- above measured shared-filesystem
throughput, where pre-building is a net loss, and below local NVMe, where it
is roughly break-even. The honest position is that this is
configuration-dependent rather than decided.

\end{document}
.
%
% Packages used are all TeX Live / Overleaf defaults: geometry, booktabs,
% amsmath, xspace, algorithm, algpseudocode, hyperref.
%
% NOT COMPILED. No LaTeX toolchain was available on the machine this was
% written on, so treat a first build as a proofreading pass.
%
% PROVENANCE RULE, inherited from the plan and enforced throughout: every
% number carries "(measured)" or "(simulated)". Simulated numbers come from
% scripts/sim_residency_v2.py over synthetic schedules with measured resource
% constants; nothing labelled simulated ran on a GPU. A retention percentage
% is never quoted without the oracle-vs-LRU gap at the same budget.

\documentclass[10pt,twocolumn]{article}
\usepackage[margin=0.75in]{geometry}
\usepackage{booktabs}
\usepackage{amsmath}
\usepackage{xspace}
\usepackage{algorithm}
\usepackage{algpseudocode}
\usepackage[hidelinks]{hyperref}

\newcommand{\sysname}{Tandem\xspace}   % placeholder; rename in one place
\newcommand{\meas}{\textsuperscript{m}}
\newcommand{\simu}{\textsuperscript{s}}

\title{\sysname: Arbitrating Model and Data Residency\\Under One Memory Budget}
\date{}

\begin{document}
\maketitle

\noindent\textbf{Notation.} \meas{} marks a measured quantity;
\simu{} marks a simulation result over measured constants.

\section{Design Overview}

An agentic HPC workflow alternates between two kinds of resource: language
models that back the planner and its tool-calling steps, and scientific data
artifacts that the tools consume. Both are expensive to make usable and both
are used repeatedly, so both are candidates for retention. \sysname is a
runtime that decides, continuously and under a single host-memory budget,
which of them to keep resident and at which level.

The design rests on three claims, each measured before it was designed
around.

\paragraph{Claim 1: making a resource usable is a chain of costs, and
residency is a prefix of that chain already paid.}
We model residency as four rungs (Table~\ref{tab:ladder}). A scheduler's job
is to choose, per resource, the highest rung it can afford to hold.

\begin{table}[t]
\centering\small
\begin{tabular}{@{}llll@{}}
\toprule
Rung & State & Cost to reach & Kind \\
\midrule
R0 & on disk & --- & --- \\
R1 & bytes in page cache & disk $\to$ RAM & movement \\
R2 & bytes in a process & RAM $\to$ process & movement \\
R3 & activated in a live consumer & parse / construct & transformation \\
\bottomrule
\end{tabular}
\caption{The residency ladder. Byte-oriented tiering systems operate on
R0$\to$R1 exclusively; that is the definition of their interface, not a
limitation of their implementation.}
\label{tab:ladder}
\end{table}

\paragraph{Claim 2: the two resource classes sit at opposite ends of that
chain, so the correct action differs in kind.}
For a language model, \textbf{74--81\%\meas{}} of a cold boot is moving
weights and engine initialisation is $\approx$4.7\%\meas{}; the cost is
movement. For a data artifact the split inverts. Decomposing a 3.32\,GB
tabulated interatomic potential with per-process \texttt{getrusage} gives
2.0\%\meas{} disk$\to$RAM, 3.9\%\meas{} RAM$\to$process, and
\textbf{93.0\%\meas{}} parse and spline construction.

The consequence is not a tuning difference. Because movement dominates a
model, the winning action is to \emph{skip} it by never falling below R2, not
to schedule it earlier: parking eliminates both the disk$\to$RAM and
RAM$\to$GPU legs, whereas warming the page cache removes only the first and
changed measured boot time by $-7.0\%$ to $+37.1\%$\meas{} across eight nodes
--- sign-unstable. Because transformation dominates a data artifact,
pre-\emph{placing} its bytes recovers at most the 2.0\% movement share; the
only winning action is to have already run the transform, i.e. to hold the
artifact at R3 inside a live consumer.

\paragraph{Claim 3: the two savings live in the same place, which is what
makes this one system rather than two.}
A model parked at R2 and a database activated at R3 are both
\emph{anonymous memory inside a live process}. Neither is a file range, so
neither can be named, evicted, or repopulated by a byte-oriented tier; and
both draw on the same host-memory allocation. Two mechanisms with
incommensurable value functions, competing for one budget, is an arbitration
problem, and it is the contribution of this work.

\subsection{Architecture}

\sysname has four components:

\begin{enumerate}\itemsep2pt
\item \textbf{M1, model parking.} Park an outgoing inference engine at R2
      rather than terminating it.
\item \textbf{M2, activated-data retention.} Hold a parsed artifact at R3 in
      a resident worker rather than re-parsing per tool call.
\item \textbf{The arbitrator.} At every residency decision point, choose the
      set of resources to hold, across both classes, under one budget.
\item \textbf{The slack stager.} Opportunistically climb rungs for resources
      predicted to be needed soon, but only when doing so evicts nothing.
\end{enumerate}

\noindent A \emph{horizon predictor} supplies the single input the arbitrator
needs --- time until each resource is next required. Section~\ref{sec:pred}
explains why it is deliberately a secondary component.

\section{Mechanisms}

\subsection{M1: model parking}

M1 uses the inference server's existing sleep facility to move an engine's
weights from device memory into host memory while keeping the process, its
CUDA context, and its allocator alive. Measured on a 32B model at tp=2, a
cold boot of \textbf{782.27\,s} becomes a \textbf{2.076\,s} wake\meas{}, with
verbatim-identical generated text, unchanged \texttt{finish\_reason}, and
post-wake throughput of 20.05--20.08\,tok/s against a 20.20\,tok/s pre-sleep
steady state --- so no cost is deferred into generation. Wake is independent
of page-cache state (evicting 18 weight shards changed it by
0.015\,s\meas{}), which is the mechanistic confirmation that the weights are
held in anonymous memory rather than re-read from a file.

Parking is not free: it costs \textbf{1.90$\times$}\meas{} the weight bytes
in host RAM (120.77\,GiB to park a 68.28\,GB model). That ratio is what makes
residency a budgeted decision rather than a default.

\emph{We do not claim sleep/wake as a contribution.} It is a feature of the
serving stack, obtainable from a flag. What \sysname contributes on this axis
is deciding \emph{when} parking is affordable and \emph{what it should
displace} --- and, critically, that at realistic allocations the answer is
frequently ``not this model.''

\subsection{M2: activated-data retention}

Agent frameworks typically execute each tool call as a fresh process. The
activated structure therefore dies at every invocation, and the artifact is
re-parsed even though the bytes are unchanged and the page cache is already
warm --- a cost invisible to any byte-oriented prefetcher by the second
parse. M2 replaces fork-per-call with a resident worker that retains the
activated structure across calls: \textbf{93.73\,s $\to$ 10.56\,s}\meas{}
(8.88$\times$) per redundant invocation, with the physical output
bit-identical (relative difference $0.000\mathrm{e}{+}00$).

The mechanism is not free of semantics, and this is where most of the
engineering lives. A live simulation instance refuses to redefine its
simulation box, and the conventional reset command destroys precisely the
tables being retained. Reaching a new problem geometry while preserving them
requires an in-place empty--resize--refill sequence. Because \sysname
rewrites commands to achieve this, \textbf{a run that completes and returns
plausible numbers is not evidence of correctness}; retention is enabled only
for command shapes that have passed a differential test executing the same
script through both the retained and the fork-per-call path and comparing
physical outputs.

\section{Arbitration}

\subsection{One budget, two classes}

Let $\mathcal{R}$ be the resident-candidate set, $g(r)$ the host memory a
resource occupies while held, $B$ the budget, and $\Delta t(r)$ the predicted
time until $r$ is next needed. Define

\begin{align}
\mathrm{full}(r) &= \mathrm{cold}(r) - \mathrm{ready}(r) \\
\mathrm{benefit}(r) &=
  \begin{cases}
    \mathrm{full}(r) & \text{retain}\\
    \min\!\left(\mathrm{full}(r), \Delta t(r)\right) & \text{stage}
  \end{cases}\\
v(r) &= \mathrm{benefit}(r)\cdot \frac{H}{\max(\Delta t(r),\,H)}
\label{eq:value}
\end{align}

\noindent where $\mathrm{cold}$ is the cost of producing the resource from
R0, $\mathrm{ready}$ the cost of resuming it from its held rung, and $H$ a
planning horizon. The arbitrator solves
$\max \sum_{r\in S} v(r)$ subject to $\sum_{r \in S} g(r) \le B$.

\subsection{Why this value function and not the obvious one}

The natural choice --- benefit per byte per second, $v = \mathrm{benefit} /
\Delta t$ --- is wrong, and wrong in a way that dimensional analysis exposes
before any experiment. Retain and stage have benefit functions of different
shape in $\Delta t$: retain's is constant, stage's saturates at
$\min(\mathrm{full}, \Delta t)$. Dividing by $\Delta t$ therefore decays
retain's score as $1/\Delta t$ while leaving stage's at a flat $1/g(r)$
whenever the window is shorter than the load --- which is every interesting
case, gaps being seconds against loads of hundreds of seconds. The resulting
density has units of $\mathrm{GB}^{-1}$, not seconds per GB. Empirically, a
stage worth 5\,s outranks and evicts a retention worth 798\,s.

We test candidate objectives against an \textbf{oracle-monotonicity
invariant}: offering staging as an additional \emph{option} to an
oracle-informed policy must never lower the result, since the policy could
always decline. The rate objective violates it by $-2.05$ to $-3.77$
points\simu{}; the pure-benefit objective is marginal; the horizon objective
of Eq.~\ref{eq:value} passes at $+3.14$ to $+4.28$\simu{}. The invariant is
necessary, not sufficient --- all arms remain myopic one-step solves --- but
it is cheap and it caught a defect that had survived in an earlier
implementation.

$H$ is a real knob and is swept rather than chosen. Smaller $H$ is better
because the pathology lives at small $\Delta t$; any $H$ above the typical
inter-need gap suffices.

\subsection{Why the problem is 0/1 and not fractional}

Eviction here is a \emph{semantic} operation, not a paging one. Discarding a
model's weights is all-or-nothing, and terminating a worker destroys the
whole activated structure; there is no interface for parking half a model or
retaining 40\% of a database, and 90\% of a model's weights serve zero
tokens. Consequently the items are indivisible and the problem is 0/1
knapsack, for which density-ranked greedy is a heuristic rather than an
optimum. This is why \sysname evaluates a subset choice rather than
single-victim eviction: single-victim greedy cannot express ``evict these two
together to admit that one.'' Algorithm~\ref{alg:arb} states the decision
procedure; note that both retain and stage candidates are scored into the
\emph{same} objective, which is what makes the two actions comparable rather
than separately prioritised.

\begin{algorithm}[t]
\small
\caption{Arbitration at a residency decision point}
\label{alg:arb}
\begin{algorithmic}[1]
\State \textbf{input:} candidates $\mathcal{R}$, budget $B$, horizon oracle $\Delta t(\cdot)$
\State $\mathcal{A} \gets \{\,r \in \mathcal{R} : \Delta t(r) \ne \textsc{Never}\,\}$
  \Comment{believed-dead excluded, not scored zero}
\For{$r \in \mathcal{A}$}
  \State $w_r \gets v(r)$ via Eq.~\ref{eq:value} with kind \textsc{retain}
\EndFor
\For{$r \in \textsc{StageCandidates}$}
  \State $w_r \gets v(r)$ with kind \textsc{stage}
  \Comment{same currency as retain}
\EndFor
\State $S \gets \textsc{Knapsack}(w, g, B)$
  \Comment{ties break toward smaller $\sum g$}
\State \textbf{evict} $\mathrm{resident} \setminus S$; \textbf{admit} $S \setminus \mathrm{resident}$
\State \Return $S$
\end{algorithmic}
\end{algorithm}

\subsection{Interaction with device slots}

A model resident on an accelerator does not charge the host-memory budget,
since its weights are in device memory. The arbitrator therefore treats
device residency as a separate, smaller capacity and only considers for
host-memory retention what has been displaced from it. This has a
counter-intuitive consequence that the evaluation confirms: adding a device
slot \emph{reduces} the value of cost-aware host retention, because the
frequently-used models are already free and what remains competing for host
memory is the less-used tail.

\section{Speculative Staging Under Slack}

\sysname stages speculatively only under a strict safety rule: \textbf{a
stage may never cause an eviction}. If the predicted-next resource does not
fit in currently unused budget, the stage is abandoned rather than funded by
displacing a resident resource.

The rule is what makes staging safe at realistic prediction quality. Because
a stage can only ever consume memory nothing else wanted, its downside is
bounded by wasted bandwidth rather than by a lost retention, and its benefit
becomes nearly independent of predictor accuracy: $+3.83$ points at 0.40
accuracy versus $+5.06$ at a perfect oracle\simu{}. \emph{The safety rule
substitutes for accuracy.} This is the inverse of the usual prefetching
design point, where value scales with precision, and it is the direct
consequence of operating in a regime where the competing use of memory is
retention rather than nothing.

\section{The Horizon Predictor, and Why It Is Secondary}
\label{sec:pred}

\subsection{From identity to horizon}

An earlier version of this system used prediction the conventional way: infer
\emph{which} resource is needed next and fetch it. That framing failed, and
its failure motivated the present design. Two signals were available --- a
central planner's plan, and learned tool-to-tool transition statistics --- and
neither delivered usable identity prediction on non-deterministic agent
traces.

Under \sysname prediction answers a different question: \emph{how long until
resource $r$ is needed again?} It is queried per already-held resource at
eviction time, not consulted to select a fetch target. This is a strictly
weaker requirement, in a way that matters:

\begin{itemize}\itemsep2pt
\item A retention decision is made \emph{at eviction time}, when the past is
      known, rather than at fetch time, when it is not.
\item A wrong retain costs only memory. A wrong fetch costs bandwidth
      \emph{and} memory \emph{and} may leave the consumer stalled anyway.
\item Identity is not required at all: every held resource is scored, so
      there is no candidate to name.
\end{itemize}

\subsection{The accuracy the design actually needs}

Because the arbitrator ranks rather than selects, horizon error degrades the
ranking gracefully instead of committing to a wrong target. Across a sweep of
predictor accuracy from 0.15 to 1.00, the deployable density-ranked policy
moves from 47.63\% to 64.86\%\simu{} while a size-blind LRU baseline is flat
at 56.23\%\simu{} by construction. Break-even against LRU falls near
\textbf{0.55} --- the middle of the accuracy band the underlying predictor
actually achieves.

We state the negative half plainly: \emph{below that break-even, cost-aware
retention is worse than LRU}, and the entire oracle ceiling over LRU is about
nine points. The predictor is secondary not because it is unimportant but
because the mechanism captures most of the available benefit without it, and
the policy's marginal contribution is real but bounded.

\subsection{One non-obvious interface requirement}

A predictor must never report ``never again.'' In a finite trace, ``no
further use'' is an artifact of the trace ending, but the arbitrator treats it
as licence to discard. A bounded-lookahead predictor should instead report
\emph{not within $H$}, which preserves a resource's eligibility when budget
is available. Making this substitution is worth roughly nine points\simu{} to
the subset-selection arm at measured accuracy, and it converts an apparent
finding --- that exact optimisation amplifies prediction error --- into an
artifact of a dead/alive binary. Table~\ref{tab:iface} states the full
interface.

\begin{table}[t]
\centering\small
\begin{tabular}{@{}p{0.30\columnwidth}p{0.60\columnwidth}@{}}
\toprule
Requirement & Why \\
\midrule
Horizon in \emph{seconds} & The budget is a rate-vs-time trade; step counts
  do not compose with load costs. \\
Queried per \emph{held} resource & The decision is which of $n$ residents to
  drop, not which one thing to fetch. \\
Calibrated ``not within $H$'' & Never $\infty$; see above. \\
Queryable \emph{at eviction time} & The decision point is displacement, not
  tool dispatch. \\
Held footprint, not file size & The budget is charged the parked/activated
  size, e.g. $1.90\times$ weights. \\
Cold and ready costs, separately & $v(\cdot)$ needs
  $\mathrm{cold}-\mathrm{ready}$, not one load estimate. \\
Per-resource confidence & Used as a weight, so a low-confidence horizon
  degrades a rank rather than forcing a choice. \\
\bottomrule
\end{tabular}
\caption{What the horizon predictor must expose.}
\label{tab:iface}
\end{table}

\section{Scope Conditions}

The design is not universally profitable, and the conditions are sharp enough
to state rather than let a reader discover.

\paragraph{It requires a stall-dominated workload.} The benefit is a
fraction of time spent waiting for resources; where compute windows are large
the same seconds are a small share of wall time. This is defensible for the
workloads measured here --- a representative trial spends 280.3\,s of
5288.5\,s\meas{} doing agent work --- but it is a real boundary, and it runs
opposite to what speculative staging wants. \textbf{The two halves of the
system prefer opposite regimes}, which we report rather than reconcile.

\paragraph{It requires contention.} If the budget accommodates every
resource, the correct policy is ``hold everything'' and no arbitration is
needed. Conversely, if the largest resources exceed the budget alone they are
not badly ranked but unrankable --- at a 256\,GB allocation the two largest
models are individually 279.0 and 276.3\,GB\meas{}, so 51.4\%\simu{} of stall
is capacity rather than replacement, and no ranking can convert it.
Quantisation, not a better policy, is what moves that mass.

\paragraph{The value of each mechanism depends on the budget, and must be
quoted with it.} At a budget where nearly everything fits, retention is worth
a great deal and an oracle recovers essentially nothing over LRU --- the win
belongs to the serving stack's flag, not to arbitration. At the production
allocation the ordering inverts, and retaining \emph{both} classes under LRU
is worse than retaining data alone by 4.73 points\simu{}, because a
size-blind policy parks a large, rarely-used model over a smaller,
frequently-used artifact; cost-aware arbitration recovers $+1.29$\simu{}.
\textbf{Any retention figure reported without the oracle-vs-LRU gap at the
same budget is not a claim about this system.}

\paragraph{Transformation cost is format-determined, and formats can be
changed.} Seconds per GB retained spans $65\times$ across
formats\meas{} and is flat to $1.00$--$1.16\times$ across a $4\times$ size
range\meas{}, which is what makes it a cacheable scheduling input rather than
an online-learning problem. But an artifact whose consumer accepts a
pre-built binary form can have its transformation cost converted away
offline. That conversion is not free: it trades a low-variance CPU cost for a
larger byte-movement cost, and for the sequence database measured here the
two balance at \textbf{151\,MB/s}\meas{} --- above measured shared-filesystem
throughput, where pre-building is a net loss, and below local NVMe, where it
is roughly break-even. The honest position is that this is
configuration-dependent rather than decided.

\end{document}
